\documentclass[tikz,border=4pt]{standalone}
\usepackage{ctex}
\usepackage{amsmath,amssymb}
\usepackage{esint}
\usetikzlibrary{calc,arrows.meta,positioning}
\begin{document}
\begin{tikzpicture}[
  font=\small,
  >=Stealth,
  box/.style={
    draw=gray!60, fill=gray!8, rounded corners=4pt,
    minimum width=4.5cm, minimum height=2.6cm,
    align=center, inner sep=4pt
  },
  arrow/.style={->, very thick, blue!50!black}
]
% ===== 中心标题 =====
\node[draw=blue!60, fill=blue!10, rounded corners=4pt,
      font=\Large\bfseries, align=center, inner sep=8pt,
      minimum width=6cm]
  (center) at (0,0) {统一原理\\
      $\displaystyle\int_{\partial M}\!\omega = \int_M d\omega$\\
      \footnotesize{(广义 Stokes 定理)}};

% ===== 四角四个定理 =====
% 左上：Newton-Leibniz (1D)
\node[box, anchor=center] (NL) at (-6, 3) {
  \textbf{Newton-Leibniz}\\
  ($n=1$)\\[3pt]
  $\displaystyle\int_a^b f'(x)\,dx=f(b)-f(a)$\\[3pt]
  \footnotesize 边界 = 两端点\\
  \footnotesize 内部 = 一维区间
};

% 右上：Green (2D plane)
\node[box, anchor=center] (GR) at (6, 3) {
  \textbf{Green 公式}\\
  ($n=2$，平面)\\[3pt]
  $\displaystyle\oint_{\partial D}\!\!P\,dx+Q\,dy=\iint_D\!\!(Q_x-P_y)\,dA$\\[3pt]
  \footnotesize 边界 = 闭曲线\\
  \footnotesize 内部 = 平面区域
};

% 左下：Stokes (2D in 3D)
\node[box, anchor=center] (ST) at (-6, -3) {
  \textbf{Stokes 公式}\\
  (3D 中曲面)\\[3pt]
  $\displaystyle\oint_{\partial\Sigma}\!\!\mathbf{F}\!\cdot\!d\mathbf{r}=\iint_\Sigma(\nabla\times\mathbf{F})\!\cdot\!d\mathbf{S}$\\[3pt]
  \footnotesize 边界 = 闭曲线\\
  \footnotesize 内部 = 3D 中曲面
};

% 右下：Gauss (3D)
\node[box, anchor=center] (GS) at (6, -3) {
  \textbf{Gauss 公式}\\
  ($n=3$)\\[3pt]
  $\displaystyle\!\!\oiint_{\partial\Omega}\!\!\!\!\mathbf{F}\!\cdot\!d\mathbf{S}=\iiint_\Omega\!\!\nabla\!\cdot\!\mathbf{F}\,dV$\\[3pt]
  \footnotesize 边界 = 闭曲面\\
  \footnotesize 内部 = 3D 区域
};

% ===== 箭头 =====
\draw[arrow] (NL.east) -- (center.north west);
\draw[arrow] (GR.west) -- (center.north east);
\draw[arrow] (ST.east) -- (center.south west);
\draw[arrow] (GS.west) -- (center.south east);

% ===== 底部说明 =====
\node[font=\small, align=center] at (0, -5.5) {
  \textbf{共同模式}：边界上的低维积分 $=$ 内部的导数积分\\
  $\partial$（边界算子） 与 $d$（外微分） 互为对偶
};
\end{tikzpicture}
\end{document}
